%% LyX 2.4.0~RC3 created this file.  For more info, see https://www.lyx.org/.
%% Do not edit unless you really know what you are doing.
\documentclass[spanish]{article}
\usepackage[T1]{fontenc}
\usepackage[utf8]{inputenc}
\usepackage{refstyle}
\usepackage{amssymb}
\usepackage{graphicx}

\makeatletter

%%%%%%%%%%%%%%%%%%%%%%%%%%%%%% LyX specific LaTeX commands.

\AtBeginDocument{\providecommand\figref[1]{\ref{fig:#1}}}
\RS@ifundefined{subsecref}
  {\newref{subsec}{name = \RSsectxt}}
  {}
\RS@ifundefined{thmref}
  {\def\RSthmtxt{theorem~}\newref{thm}{name = \RSthmtxt}}
  {}
\RS@ifundefined{lemref}
  {\def\RSlemtxt{lemma~}\newref{lem}{name = \RSlemtxt}}
  {}


\makeatother

\usepackage{babel}
\addto\shorthandsspanish{\spanishdeactivate{~<>}}
\deactivatequoting

\begin{document}
For the section of bifurcations...

\begin{figure}[h]
\begin{centering}
\includegraphics[width=1\textwidth]{Bifurcations_plot}
\par\end{centering}
\caption{Some bifurcations and $\mathbb{M}_{2,z_{0}\protect\neq0}$.}\label{fig:M2BifParams}
\end{figure}


\section*{Appendix: Numerical approximations and drawings}

\begin{figure}[h]
\begin{centering}
\includegraphics[width=0.75\textwidth]{Mandelbrot_ref1}
\par\end{centering}
\caption{$\mathbb{M}$, the Mandelbrot set, drawn in black. The other colors
are determined by the escape time of the iterates $z_{n}=q_{c}^{n}(0)$,
considering that $z_{n}$ escapes when $|z|^{2}>4$.}\label{fig:TheMandelbrot}
\end{figure}

In order to better understand this section, let us recall that the
Mandelbrot set is
\[
\mathbb{M}=\left\{ c\in\mathbb{C}|\,\,o(0,q_{c})\,\,\mathrm{is\,bounded}\right\} =\left\{ c\in\mathbb{C}|\,\,J(q_{c})\,\,\mathrm{is\,connected}\right\} ,
\]
 where $q_{c}(z)=z^{2}+c$. The point $0$ is the critical point of
each function $q_{c}$. The classical algorithm of coloring by escape
time is commonly used to draw the Mandelbrot set (see \figref{TheMandelbrot}).

The family $f_{\lambda,\mu}$ does not have a unique set of critical
points for all pairs $(\lambda,\mu)$; there are infinitely many critical
points for each $(\lambda,\mu)$, and none of them can be written
explicitly in a formula. Thus, the drawing of an analogous Mandelbrot
set cannot directly use the classical drawing algorithms.

First, it is necessary to calculate the critical points. Let us recall
that the critical points $c_{k}=c_{k}(\lambda,\mu)$ of $f_{\lambda,\mu}$
corresponds to the intersections of the curves

\[
y_{\pm}(x)=\pm\sqrt{\left|\frac{\mu}{\lambda}\right|\frac{1}{e^{x}}-x^{2}},
\]
\[
x(y)=-\frac{y}{\tan(\frac{y+\pi}{2})}.
\]

In the following items, we describe the algorithms to approximate
numerically the critical points $c_{k}=c_{k}(\lambda,\mu)$, for $k\in\mathbb{N}-\left\{ 0\right\} $,
given a parameter pair $(\lambda,\mu)\in\mathbb{R}^{-}\times\mathbb{R}^{+}$. 
\begin{enumerate}
\item Approximation of $c_{1}$, the nearest critical point to $0$ with
$Re(c_{1})>0$:
\begin{enumerate}
\item We know that $c_{1}$ is the intersection of the curves $y_{+}(x)$
and $x(y)$, with $y\in(0,\pi)$.
\item $c_{1}$ is approximated using a bisection search algorithm over $y$
with initial values $y_{0,\min}=\delta=0.001$ and $y_{0,\max}=\pi-\delta$,
stopping when $\left|\frac{\mu}{\lambda}\right|\frac{1}{e^{x_{n}}}-x_{n}^{2}>0$
and $|f'_{\lambda,\mu}(x_{n}+y_{+}(x_{n})i)|^{2}<\varepsilon=10^{-8}$,
where $x_{n}=x(y_{n})$.
\end{enumerate}
\item Approximation of $c_{k}$, for $k>1$.
\begin{enumerate}
\item First, using the Newton-Raphson algorithm, we approximate the root
$x_{L}<0$ of $g(x)=\left|\frac{\mu}{\lambda}\right|\frac{1}{e^{x}}-x^{2}$
such that for all $x\in(-\infty,x_{L})$ we have $g(x)>0$. Note that
$y_{+}(x)=\sqrt{g(x)}$ and $y_{+}(x_{L})=0$.
\item We know that $c_{k}\in(-\infty,x_{L})\times((2k-3)\pi,(2k-1)\pi)$.
Then, $c_{k}$ is approximated using a compass search algorithm with
initial value $z_{0}=x_{L}+2(k-1)\pi i$, stopping when $|f'_{\lambda,\mu}(z_{n})|^{2}<\varepsilon=10^{-8}$.
\end{enumerate}
\end{enumerate}
%
\begin{figure}[h]
\begin{centering}
\includegraphics[width=1\textwidth]{BDomainPlus_plot_light}
\par\end{centering}
\caption{The set $\mathbb{M}_{1}$ drawn in black. The other colors are determined
by the escape time of the iterates $z_{n}=f_{\lambda,\mu}^{n}(c_{1})$,
considering that $z_{n}$ escapes when $Re(z)<-16$.}\label{fig:M1}
\end{figure}

Now, it can be defined the analogue of the Mandelbrot set for the
family $f_{\lambda,\mu}$ (see \figref{M1}):
\[
\ensuremath{\mathbb{M}_{1}=\{(\lambda,\mu)\in\mathbb{R}^{-}\times\mathbb{R}^{+}|\,\,o(c_{1}(\lambda,\mu),f_{\lambda,\mu})\,\,\mathrm{is\,bounded}\}.}
\]

\begin{figure}[h]
\begin{centering}
\includegraphics[width=0.5\textwidth]{Mandelbrot_ref3}\includegraphics[width=0.5\textwidth]{Mandelbrot_ref2}
\par\end{centering}
\caption{At left, the bulbs of period $1$ and $2$ of the Mandelbrot set
$\mathbb{M}$. At right, the bulbs of period $1$ and $2$, and the
complement of $\mathbb{M}$ considered as a bulb of period $1$ (because
$\infty$ is a super-attracting fixed point for all $q_{c}$).}\label{fig:TheMandelbrot2}
\end{figure}

The bulbs of period $N$ of the Mandelbrot set are the open subsets
$U\subset\mathbb{M}$ such that for every $c\in U$ the corresponding
function $q_{c}$ has a finite attracting periodic orbit of period
$N$, and it is known that the critical point $0$ belongs to the
immediate basin of attraction of the orbit. The algorithm to draw
these bulbs of period $1$ and $2$ use that $\mathbb{C}$ is a complete
metric space and the convergence of Cauchy sequences: the iteration
is stopped when $|q_{c}^{n}(0)-q_{c}^{n+2}(0)|<\varepsilon$ and the
points are coloring according to $n$. See \figref{TheMandelbrot2}.

\begin{figure}[h]
\begin{centering}
\includegraphics[width=0.75\textwidth]{BulbsPer2_plot}
\par\end{centering}
\caption{The bulbs of period $1$ and $2$, and the bulb associated to the
Baker domain $F_{0}$ of the set $\mathbb{M}_{2}$, are drawn with
different colors. Note the line $\mathcal{L}_{-1}$ in light blue,
corresponding to parameters for which $f_{\lambda,\mu}$ has a fixed
parabolic point with multiplier equal to $-1$.}\label{fig:M2}
\end{figure}

Let $z_{0}=z_{0}(\lambda,\mu)\in\mathbb{C}$ be the convergence point
of the sequence $f_{\lambda,\mu}^{2n}(c_{1})$, such that $|(f_{\lambda,\mu}^{2n})'(z_{0})|<1$
(that is, $z_{0}$ is attractor) or $z_{0}=0$ (the pole). We define
\[
\ensuremath{\mathbb{M}_{2}=\{(\lambda,\mu)\in\mathbb{R}^{-}\times\mathbb{R}^{+}|\,\,f_{\lambda,\mu}^{2n}(c_{1})\rightarrow z_{0}\}},
\]
which represents the bulbs of period $1$ and $2$, and also the ``Baker
bulb'' corresponding to the parameters for which $f_{\lambda,\mu}^{2n}(c_{1})\rightarrow0$.
See \figref{M2}.

\begin{figure}[h]
\begin{centering}
\includegraphics[width=0.5\textwidth]{BulbsPer2Finite_plot}\includegraphics[width=0.5\textwidth]{BulbBaker_plot}
\par\end{centering}
\caption{At left $\mathbb{M}_{2,z_{0}\protect\neq0}$, the bulbs of period
$1$ and $2$. At right, the ``Baker bulb'' $\mathbb{M}_{2,z_{0}=0}$.}\label{fig:M2Bulbs}
\end{figure}

Naturally, one can differentiate the set of bulbs of period $1$ and
$2$, and the ``Baker bulb'' set:

\[
\ensuremath{\mathbb{M}_{2,z_{0}\neq0}=\{(\lambda,\mu)\in\mathbb{R}^{-}\times\mathbb{R}^{+}|\,\,f_{\lambda,\mu}^{2n}(c_{1})\rightarrow z_{0}\neq0\}\subset\mathcal{M}_{2},}
\]

\[
\ensuremath{\mathbb{M}_{2,z_{0}=0}=\{(\lambda,\mu)\in\mathbb{R}^{-}\times\mathbb{R}^{+}|\,\,f_{\lambda,\mu}^{2n}(c_{1})\rightarrow0\}\subset\mathcal{M}_{2}.}
\]

See \figref{M2Bulbs} for drawings of $\mathbb{M}_{2,z_{0}\neq0}$
and $\mathbb{M}_{2,z_{0}=0}$. 

\section*{Appendix: Conjectures}

Let us recall that
\[
\ensuremath{\mathbb{B}=\{(\lambda,\mu)\in\mathbb{R}^{-}\times\mathbb{R}^{+}|\,\,F_{0}\cap Crit(f_{\lambda,\mu})\neq\emptyset\}},
\]
where $F_{0}=F_{0}(\lambda,\mu)$ is the Baker domain of period $2$
such that $0\in\partial F_{0}$, and $f_{\lambda,\mu}(F_{0})=F_{\infty}$.

\paragraph*{Conjecture 1.}\label{par:Conjecture1.}

If $(\lambda,\mu)\in\mathcal{\mathbb{B}}$, then $c_{1},c_{-1}\in F_{0}$,
and $c_{k},c_{-k}\in F_{\infty}$ for all $k>1$.\\

According to Conjecture 1, $\ensuremath{\mathbb{B}}$ corresponds
to the component of $\mathbb{C}-\mathbb{M}_{1}$ containing the parameter
pair $(-20,1/4)$, drawn in white and other colors (but not black)
in \figref{M1}. Then we have

\paragraph*{Corollary of Conjecture 1.
\[
\mathbb{B}\subset\mathcal{\mathbb{M}}_{2,z_{0}=0}.
\]
}

\begin{figure}[h]
\begin{centering}
\includegraphics[width=0.75\textwidth]{Conjecture_plot}
\par\end{centering}
\caption{Parameter pairs in $\mathbb{M}_{2,z_{0}=0}$, but outside of $\mathbb{B}$.}\label{fig:M2Conjecture}
\end{figure}

\begin{figure}[h]
\begin{centering}
\includegraphics[width=0.75\textwidth]{DynPlane_l-1698_m5}
\par\end{centering}
\caption{$(\lambda,\mu)=(-5e/8,5)$. $f_{\lambda,\mu}^{n}(c_{1})\protect\notin F_{0}$
for $n=0,1,2$; $f_{\lambda,\mu}^{3}(c_{1})\in F_{0}$.}\label{fig:DynConjecture1}
\end{figure}

\begin{figure}[h]
\begin{centering}
\includegraphics[width=0.75\textwidth]{DynPlane_l-3058_m475}
\par\end{centering}
\caption{$(\lambda,\mu)=(-9e/8,4.75)$. $f_{\lambda,\mu}^{n}(c_{1})\protect\notin F_{0}$
for $n=0,1,2,3,4$; $f_{\lambda,\mu}^{5}(c_{1})\in F_{0}$.}\label{fig:DynConjecture2}
\end{figure}

\begin{figure}[h]
\begin{centering}
\includegraphics[width=0.75\textwidth]{DynPlane_l-3624_m55}
\par\end{centering}
\caption{$(\lambda,\mu)=(-4e/3,5.5)$. $f_{\lambda,\mu}^{n}(c_{1})\protect\notin F_{0}$
for $n=0,1,2,3$; $f_{\lambda,\mu}^{4}(c_{1})\in F_{0}$.}\label{fig:DynConjecture3}
\end{figure}


\paragraph*{Conjecture 2.}

For all $(\lambda,\mu)\in\mathbb{R}^{-}\times\mathbb{R}^{+}$ and
for any $k\in\mathbb{Z}-\left\{ -1,0,1\right\} $, $c_{k}\in F_{\infty}$.\\


\paragraph*{Corollary of Conjecture 2.}

For any $k\in\mathbb{Z}-\left\{ -1,0,1\right\} $
\[
\mathbb{M}_{2,z_{0}=0}=\left\{ (\lambda,\mu)\in\mathbb{R}^{-}\times\mathbb{R}^{+}|\,\,f_{\lambda,\mu}^{2n}(c_{k})\rightarrow\infty\right\} 
\]
\[
=\left\{ (\lambda,\mu)\in\mathbb{R}^{-}\times\mathbb{R}^{+}|\,\,f_{\lambda,\mu}^{2n+1}(c_{k})\rightarrow0\right\} .
\]


\paragraph*{Conjecture 3.}

$\mathbb{B}$ is simply connected and $\mathcal{\mathbb{M}}_{2,z_{0}=0}$
is not connected.\\

In \ref{fig:M1}, \ref{fig:M2}, \ref{fig:M2Bulbs}, and \ref{fig:M2Conjecture},
we observe a subset $Q$ with the familiar shape of the Mandelbrot
set. We state the following 

\paragraph*{Conjecture 4.}

$f_{\lambda,\mu}$ is a family of quadratic-like maps in a set of
parameters $Q\subset[-4.5,-3.5]\times[4,5]$.\\


\paragraph*{Conjecture 5.}

There are infinite many subsets of parameters for which $f_{\lambda,\mu}$
is a family of quadratic-like maps.
\end{document}
